% résumé


\setstretch{1.5}

Le résumé se trouve possiblement à trois endroits:
\begin{enumerate}
	\item  juste après la page titre
	\item  sur la dernière page
	\item  sur la quatrième page de garde (dos de la couverture du rapport)
\end{enumerate}


Il est très souvent écrit en anglais et accompagné de \textbf{cing mots clés}
qui permettent ensuite un indexage plus facile dans les bibliothèques ou sur internet.

Le résumé  est écrit très souvent en anglais, et 
peut être  accompagné d'une version française.

Un résumé ne dépasse jamais une page, soit de 15 à 20 lignes environ.

Enfin on peut choisir une taille de fonte légèrement plus petite que celle
du document, ce qui permet d'allonger le texte ou de mettre les deux versions
anglaise et française sur une seule page.


Le résumé rappelle les éléments suivants:

\begin{itemize}
\item le contexte, la motivation et les objectifs du projet
\item les méthodologies adoptées pour résoudre le problème posé
\item le ou les principaux résultats originaux incluant les analyses
\item la ou les principales conclusions
\item une perspective intéressante
\end{itemize}


\setstretch{1.}
 